\section{Cómo la computadora «ve» las imágenes}

\subsection{Representación numérica de las imágenes}

Para la computadora, una imagen no es más que números. Una imagen en escala de grises es un arreglo bidimensional en el que cada elemento (píxel) es un entero en el rango $[0, 255]$ que representa el brillo de ese píxel. Una imagen a color, en cambio, es un arreglo tridimensional de tamaño $H \times W \times 3$, donde la tercera dimensión corresponde a los tres canales de color RGB.

\begin{figure}[H]
\centering
\includegraphics[width=0.9\textwidth]{slides-images/slide-12.jpg}
\caption{La imagen es números: una imagen en escala de grises es una matriz bidimensional, y una imagen a color (por ejemplo, de 1080$\times$1080) es un apilamiento de tres matrices de este tipo (los tres canales RGB).\protect\footnotemark}
\end{figure}
\footnotetext{Slides, página 12.}

\begin{importantbox}{Representación matemática de las imágenes}
\begin{itemize}
    \item \textbf{Imagen en escala de grises}: $\mathbf{I} \in \mathbb{R}^{H \times W}$, cada valor de píxel $I_{i,j} \in [0, 255]$
    \item \textbf{Imagen a color}: $\mathbf{I} \in \mathbb{R}^{H \times W \times 3}$, los tres canales corresponden respectivamente a rojo (R), verde (G) y azul (B)
    \item Por ejemplo, una imagen RGB de $1080 \times 1080$ contiene $1080 \times 1080 \times 3 = 3,499,200$ números
\end{itemize}
\end{importantbox}

Amini señala que, desde el punto de vista de la representación, las imágenes en realidad son más fáciles de procesar que el lenguaje, porque las imágenes ya son de manera natural arreglos de números, mientras que el lenguaje requiere primero una conversión, como el embedding de palabras.

\subsection{Dos tipos básicos de tareas en visión por computadora}

Los dos tipos de tareas más básicos en visión por computadora son la \textbf{regresión} (Regression) y la \textbf{clasificación} (Classification):

\begin{figure}[H]
\centering
\includegraphics[width=0.85\textwidth]{slides-images/slide-13.jpg}
\caption{Tareas básicas de la visión por computadora: clasificación (produce una etiqueta de categoría discreta y su probabilidad) y regresión (produce un valor continuo).\protect\footnotemark}
\end{figure}
\footnotetext{Slides, página 13.}

\begin{itemize}
    \item \textbf{Regresión}: la variable de salida toma valores continuos, por ejemplo, predecir el ángulo de giro del volante.
    \item \textbf{Clasificación}: la variable de salida toma etiquetas de categoría discretas, por ejemplo, determinar qué presidente estadounidense es la persona en la imagen. Una tarea de clasificación puede producir una distribución de probabilidad sobre cada categoría.
\end{itemize}

Sin importar de qué tarea se trate, el problema central siempre es la \textbf{detección de características} (Feature Detection) — entender qué características distinguen a las diferentes categorías entre sí.

\subsection{Resumen de la sección}

La computadora representa las imágenes mediante matrices de valores de píxeles. Una imagen a color es un tensor tridimensional (alto $\times$ ancho $\times$ número de canales). Los dos tipos básicos de tareas en visión por computadora —clasificación y regresión— dependen de la capacidad de extraer características de manera efectiva a partir de la imagen.

\newpage
